% Generated draft: requires mathematical and source review.
\documentclass[11pt]{article}
\usepackage[margin=25mm]{geometry}
\usepackage{fontspec}
\setmainfont{Latin Modern Roman}
\usepackage{amsmath,amssymb,mathtools,unicode-math}
\setmathfont{Latin Modern Math}
\usepackage{enumitem,needspace,xurl,hyperref}
\setlength{\parindent}{0pt}
\setlength{\parskip}{5pt}
\setlength{\emergencystretch}{3em}
\newcommand{\N}{\mathbb{N}}
\newcommand{\Z}{\mathbb{Z}}
\newcommand{\Q}{\mathbb{Q}}
\newcommand{\R}{\mathbb{R}}
\newcommand{\C}{\mathbb{C}}
\newcommand{\F}{\mathbb{F}}
\newcommand{\A}{\mathbb{A}}
\newcommand{\PP}{\mathbb{P}}
\newcommand{\E}{\mathbb{E}}
\newcommand{\eps}{\varepsilon}
\newcommand{\dd}{\,\mathrm{d}}
\begin{document}
\section*{UnsolvedMath ID: 800820}
\textbf{Source problem code:} catalogue entry
\subsection{Definitions and mathematical statement}
Let \(a>0\) and let \(M=\mathbb T\times(0,a)\) be a periodic channel, with coordinates \((x_1,x_2)\); the period of \(\mathbb T\) is not specified. Consider a plane-parallel velocity field
\[
\bar u(x_1,x_2)=(v(x_2),0),\qquad v'(x_2)>0,\quad v''(x_2)>0\quad(0<x_2<a).
\]
Let \(u_0\) be an initial velocity field for an ideal fluid in \(M\), assumed to be \(C^1\) close to \(\bar u\). Set \(\omega_0=\operatorname{curl}u_0\), and assume that this scalar vorticity has a gradient satisfying
\[
\nabla\omega_0(x)\ne0\qquad\text{for every }x\in M.
\]
Suppose that, for some value \(c\), one of the level lines of \(\omega_0=c\) has a feature called a ``kink'' in the source.

The requested assertion is that this kink does not disappear as the ideal fluid evolves in time. In particular, the intended persistence concerns the initially specified kink, rather than merely the existence of some kink somewhere at each later time.

This remains a qualified statement: the supplied record does not define a kink, specify how to identify its continuation, or give a mathematical criterion for its disappearance. It also leaves unspecified the required size and precise norm of the \(C^1\) perturbation, the evolution equations and boundary conditions, the solution regularity, and the time interval of the assertion. These choices are needed before the persistence claim can be expressed as a fully precise quantified proposition. If \(\omega_0\) is continuously differentiable, the nonvanishing-gradient hypothesis makes its level lines locally \(C^1\); consequently, interpreting ``kink'' as a literal corner would require resolving an apparent incompatibility with that regularity.
\subsection{Short English statement}
For an ideal fluid in a periodic channel, start near a parallel shear whose profile has strictly positive first and second derivatives. Assume that the initial vorticity has no critical points and that one vorticity level line has a ``kink.'' The source asks for proof that this kink persists during the evolution. The meaning of ``kink'' and several analytic details are not supplied, so this is a statement-only draft requiring clarification and human mathematical review.
\subsection{Sources}
\begin{itemize}
\item Local UnsolvedMath registry, permanent ID \(800820\), exact source problem code \texttt{catalogue entry}, title \emph{Geometric Hydrodynamics in Open Problems --- Problem 19}. The supplied registry labels the problem open; this label is not an independently verified mathematical or literature assessment.
\item Khesin, Misiolek, and Modin, \emph{Geometric Hydrodynamics in Open Problems} (2022), Problem 19, as attributed in the supplied record. Supplied source URL: \url{https://arxiv.org/abs/2205.01143}. The record reports access on 2026-10-16 and extraction by \texttt{arxiv\_latex}; these are supplied provenance details, not a claim of independent source access or verification.
\item Supplied upstream revision: \texttt{d371ee917ac9647fbb06cd6fc529deec6f3dda82}. The local record was updated on 2026-10-16.
\end{itemize}
\end{document}
